% === Introduction and motivation ===

This document describes the algorithms, data structures, and implementation
of \FractalShark{}, a high-performance, GPU-accelerated Mandelbrot set
renderer designed for extreme deep-zoom exploration.  \FractalShark{}
combines perturbation theory with linear approximation, NTT-based
high-precision GPU arithmetic, and optimized reference-orbit computation to
enable interactive navigation at zoom depths exceeding $10^{10000}$.  Its
Feature Finder automatically locates high-period periodic points, enabling
targeted exploration of deeply embedded structures without manual search.

The project has evolved through extensive experimentation with numerical
precision, memory management, and GPU kernel design.  As a result, this
document serves a dual purpose: it is both \emph{technical documentation}
of the algorithms and architecture underlying \FractalShark{}, and a
\emph{research notebook} recording the design rationale, tradeoffs, and
lessons learned during development.  Where relevant, we note alternative
approaches that were considered or attempted, and explain why the current
design was chosen.

The presentation proceeds from high-level goals and mathematical background
through the individual subsystems---rendering kernels, perturbation
pipelines, reference-orbit computation, GPU arithmetic, and the Feature
Finder---and concludes with infrastructure details and development history.

\section{\FractalShark{} Overview}
\label{sec:goal}

\begin{figure}[!htbp]
  \centering
  \includegraphics[width=0.9\linewidth]{mandelbrot.png}
  \caption[The Mandelbrot set]{The Mandelbrot set.}
  \label{fig:mandelbrot}
\end{figure}

\FractalShark{} is an interactive Mandelbrot set renderer (\cref{fig:mandelbrot}) focused on
\textbf{extreme deep-zoom exploration}, numerical correctness, and
algorithmic experimentation.  It supports zoom depths ranging from
conventional views to magnifications exceeding $10^{10000}$, while
exposing internal rendering choices and precision tradeoffs to the user.

\paragraph{Interactive navigation and view control.}
\FractalShark{} provides direct mouse-driven navigation, including centering
the view at an arbitrary point, zooming in and out at the cursor location,
stepping backward through navigation history, and invoking automatic zoom
modes.  Window management options allow toggling between windowed and
full-screen modes, including square-aspect rendering for precise analysis.

\paragraph{Built-in views for demonstration and validation.}
The application includes a set of built-in views covering a wide
range of numerical regimes: GPU precision limits, high-period locations,
known hard points, historical bug cases, regression tests, and ultra-deep
zoom demonstrations.  These views serve both as showcases of capability and
as repeatable test cases for performance and correctness.

\paragraph{Explicit algorithm selection and transparency.}
A central design goal of \FractalShark{} is to make rendering algorithms
explicit and selectable.  Users may rely on an automatic algorithm selector
or manually choose among many rendering paths (see \cref{subsec:ref-orbit-mt-tradeoff} for automatic selection), including:
\begin{itemize}
  \item GPU-based low-zoom renderers with selectable iteration precision and
        multiple numeric formats (single-, dual-, and quad-limb 32- and 64-bit
        variants, as well as deeper representations discussed subsequently);
        see \cref{sec:mandel-base-kernels}.
  \item Scaled perturbation algorithms for extending low-precision arithmetic
        to deeper zooms (\cref{sec:perturbation-concept}).
  \item Bilinear approximation (BLA v1) perturbation paths (\cref{sec:bilinear-approx}).
  \item Linear approximation (LA~v2) pipelines, including full rendering,
        linear-approximation-only, and perturbation-only variants
        (\cref{sec:la-v2-perturb}).
  \item Reference-compression–aware algorithms, including Imagina-compatible
        ``max'' compression formats (\cref{subsec:ref-orbit-compression-reuse}).
  \item CPU-only algorithms for very high precision, verification, and
        fallback scenarios.
\end{itemize}

Many of these modes are exposed explicitly for testing and comparison, and
not all are intended as default or production-quality paths.

\paragraph{Linear approximation configuration.}
Linear approximation behavior can be configured for multithreaded or
single-threaded execution and tuned via presets that prioritize accuracy,
performance, or memory usage.  These controls reflect ongoing development
and experimentation with LA parameter tradeoffs.

\paragraph{Image quality and coloring.}
\FractalShark{} supports GPU antialiasing at multiple sample levels and a
flexible palette system (\cref{subsec:antialiasing,subsec:coloring}).  Palettes may be selected from predefined themes,
generated randomly, and rendered at configurable color depths.  Some palette
features are intentionally limited or disabled where they are known to be
incomplete.

\paragraph{Iteration limits and precision control.}
Users can dynamically adjust iteration limits, switch between 32-bit and
64-bit iteration counters, and modify shallow-zoom iteration precision.
These controls are primarily relevant for low-zoom or testing scenarios and
are not universally applicable to all rendering algorithms.

\paragraph{Perturbation and reference-orbit management.}
\FractalShark{} places strong emphasis on reference-orbit reuse and
perturbation-based rendering (\cref{sec:ref-orbit-calc}).  It supports multiple perturbation strategies
(single-threaded, multithreaded, periodicity-assisted (\cref{sec:hp-periodicity}), and experimental GPU
variants), along with tools to clear, inspect, save, reload, and reuse
reference orbits.  Automatic and manual reference-orbit persistence is
supported via file-backed storage (\cref{sec:disk_backed_growable_vectors}).

\paragraph{Memory management and scalability.}
To enable extreme zoom depths without exhausting system memory, \FractalShark{}
employs file-backed storage, optional reference compression, configurable
memory limits, and automatic cleanup policies.  These mechanisms are closely
tied to recent allocator and reference-orbit refactors discussed in the
project development history.

\paragraph{Diagnostics, benchmarking, and testing.}
The application includes facilities for displaying detailed rendering
parameters, running repeatable benchmarks, executing regression tests, and
comparing reference orbits (including Imagina-compatible formats).  These
features reflect \FractalShark{}'s dual role as both a visualization tool and a
development platform.

\paragraph{Data export and interoperability.}
\FractalShark{} can save rendered images, high-resolution bitmaps, iteration
counts, and reference orbits in multiple text and compressed formats.  It
supports loading external reference orbits and matching Imagina-compatible
data for cross-tool validation.

\medskip
\noindent
Overall, \FractalShark{} functions as both a \textbf{deep-zoom Mandelbrot
explorer} and a \textbf{research and experimentation platform} for
high-precision fractal rendering, prioritizing transparency, correctness, and
performance.  It is not intended as a polished end-user application, but rather
as an experimental platform for exploring and validating advanced fractal
rendering algorithms.

\section{The Mandelbrot set and escape-time rendering}
\label{sec:mandelbrot-intro}

The Mandelbrot set \(M\subset\mathbb{C}\) \cite{Mandelbrot1980Fractal} is defined as the set of complex
parameters \(c\) for which the orbit of
\begin{equation}
z_{n+1} = z_n^2 + c,\qquad z_0 = 0
\end{equation}
remains bounded. A common rendering method is \emph{escape-time}: iterate until
the orbit magnitude exceeds a bailout threshold, or until a maximum iteration
count is reached.

A standard bailout uses \(|z_n|^2 \ge 4\), since \(|z|>2\) implies divergence:
\begin{equation}
|z_n|^2 = \Re(z_n)^2 + \Im(z_n)^2 \ge 4 \quad\Rightarrow\quad \text{escaped.}
\end{equation}
For each pixel, the stored value is the smallest \(n\) (or a bounded proxy) at
which escape occurs, or the maximum iteration limit if escape never occurs.


\subsection*{System architecture}
\label{sec:system-architecture}

At a high level, \FractalShark{}'s rendering pipeline transforms a view
specification (centre, zoom, iteration limit) into a coloured image through
a sequence of stages.  Most stages are implemented as CUDA kernels; the CPU
orchestrates data flow between them and handles high-precision coordinate
arithmetic.

\begin{verbatim}
  View parameters (centre, zoom, iterations)
          |
          v
  Coordinate mapping            [PointZoomBBConverter]
   (high-precision -> per-pixel deltas)
          |
          v
  Reference orbit computation   [RefOrbitCalc]
   (CPU: MPIR single/multi-thread,
    GPU: HpSharkFloat NTT pipeline)
          |
          v
  Linear approximation setup    [LAReference]
   (period detection, LA/BLA coefficient tables)
          |
          v
  Per-pixel perturbation (GPU)  [mandel_1xHDR_float_perturb_lav2]
   (LA stage traversal -> perturbation tail -> escape test)
          |
          v
  Antialiasing + coloring (GPU) [antialiasing_kernel]
   (supersample averaging, palette lookup)
          |
          v
  Display (OpenGL)              [RenderThreadPool GL consumer]
\end{verbatim}

The first two stages---coordinate mapping and reference orbit
computation---dominate wall-clock time at extreme zoom depths, because
the required arithmetic precision grows with the logarithm of the zoom
factor.  Once the reference orbit is available, per-pixel perturbation
on the GPU is comparatively fast: each pixel evolves only a low-precision
delta relative to the shared reference
(\cref{sec:perturbation-concept}), with linear approximation
(\cref{sec:la-v2-perturb}) skipping the majority of iterations.

Two rendering paths exist (\cref{subsec:two-render-paths}): the
\emph{render pool path}, where worker threads each acquire a private
iteration buffer and produce frames for the OpenGL consumer; and the
\emph{direct path}, used by tests and save operations, which writes
results into the shared \code{m\_CurIters} buffer.  The render thread
pool (\cref{sec:render-thread-pool}) manages a command queue that
serialises UI-driven mutations with rendering work, discarding stale
renders when the user navigates faster than the GPU can complete them.

\subsection*{Document roadmap}

The remainder of this document is organized as follows:
\begin{itemize}
  \item \textbf{Direct rendering kernels and numeric types.}
        The GPU escape-time kernels and the fixed-point and floating-point
        numeric representations they operate on
        (\cref{sec:mandel-base-kernels,sec:hdr-float,sec:extended-precision-types}).
  \item \textbf{Perturbation theory and approximation-based acceleration.}
        Scaled perturbation, bilinear approximation (BLA), and the LA~v2
        pipeline that together enable deep-zoom rendering from a single
        reference orbit
        (\cref{sec:perturbation-concept,sec:bilinear-approx,sec:la-v2-perturb}).
  \item \textbf{Reference orbit computation.}
        CPU and GPU backends for computing the high-precision reference orbits
        required by perturbation methods (\cref{sec:ref-orbit-calc}).
  \item \textbf{GPU high-precision arithmetic.}
        NTT-based multiplication, parallel carry propagation, and the
        memory layout underlying GPU-side arbitrary-precision computation
        (\cref{subsec:ref-orbit-gpu,sec:hp-add,sec:ntt-multiply}).
  \item \textbf{Feature Finder.}
        Periodic-point detection via Newton's method, period scanning, and
        automatic deep-zoom target selection (\cref{sec:feature-finder}).
  \item \textbf{Rendering pipeline and memory infrastructure.}
        Antialiasing, coloring, file-backed growable vectors, and the overall
        render orchestration (\cref{sec:misc-details,sec:render-thread-pool,sec:disk_backed_growable_vectors}).
  \item \textbf{Appendices.}
        A user guide covering the menu interface and keyboard shortcuts, and
        a development history recording implementation milestones and design
        decisions.
\end{itemize}


%% ===================================================================
\section{Template design}
\label{sec:template-design}
%% ===================================================================

\FractalShark{} makes pervasive use of C++ templates to select numeric
types, iteration widths, and algorithmic variants at compile time rather
than at run time.  Two considerations motivate this choice.  First,
the inner loops of both reference-orbit computation and per-pixel
perturbation execute billions of arithmetic operations; virtual dispatch
or run-time branch mispredictions on every iteration would impose
measurable overhead.  Second, encoding precision contracts in the type
system allows the compiler to verify that high-precision reference data
is never silently truncated when it is handed to a lower-precision
perturbation kernel, catching a class of errors that would otherwise
manifest only as subtle rendering artifacts.

The template parameters fall into four largely orthogonal dimensions,
described in the following subsections.


\subsection{Iteration-counter width (\code{IterType})}
\label{subsec:template-itertype}

Every orbit and perturbation result is indexed by an iteration counter
whose type is either \code{uint32\_t} or \code{uint64\_t}.  A 32-bit
counter suffices for the vast majority of renders, supporting up to
$2^{32}-1 \approx 4.3\times10^{9}$ iterations.  Certain ultra-deep
locations, however, require iteration limits that exceed this range.
Switching to a 64-bit counter doubles the per-element storage cost for
every orbit and iteration buffer, so \FractalShark{} treats the counter
width as a compile-time choice rather than unconditionally paying the
wider cost.  The \code{IterType} parameter threads through
\code{PerturbationResults}, \code{LAReference}, and the rendering
kernels, ensuring that all components agree on the iteration index
representation.


\subsection{Reference-orbit type (\code{T})}
\label{subsec:template-t}

The parameter~\code{T} selects the numeric type used to store the
reference orbit $\{z_n^\star\}$ and, by extension, the type in which the
high-precision reference computation is performed.  It determines both
the precision available during orbit iteration and the storage
representation uploaded to the GPU for use by perturbation kernels.
The principal choices are:

\begin{itemize}
  \item \code{float} and \code{double}---IEEE~754 single and double
        precision, used at shallow zoom depths where the full mantissa
        suffices to represent the orbit accurately.
  \item \code{CudaDblflt<MattDblflt>}---a double-single representation
        that pairs two 32-bit floats to obtain roughly double the
        mantissa bits of \code{float} alone, with natively vectorisable
        GPU arithmetic (\cref{sec:extended-precision-types}).
  \item \code{HDRFloat<double>} and \code{HDRFloat<float>}---a
        high-dynamic-range wrapper that adds an explicit exponent to the
        mantissa, preventing underflow in the deep-zoom regime where
        orbit values span hundreds of orders of magnitude
        (\cref{sec:hdr-float}).
  \item \code{HDRFloat<CudaDblflt<MattDblflt>>}---the composition of
        the HDR exponent with the double-single mantissa, providing
        both extended range and extended precision.
  \item \code{HpSharkFloat<SharkFloatParams>}---the GPU-resident
        arbitrary-precision type used for reference-orbit computation
        at extreme zoom depths (\cref{subsec:ref-orbit-gpu}).
\end{itemize}

\noindent
The type~\code{T} appears as the second template parameter of
\code{PerturbationResults<IterType,\allowbreak T,\allowbreak PExtras>}
and as the orbit-storage type inside \code{RefOrbitCalc}.


\subsection{Perturbation subtype (\code{SubType})}
\label{subsec:template-subtype}

While~\code{T} governs the precision of the shared reference orbit,
per-pixel perturbation operates at a lower precision because each
pixel evolves only a small delta $\Delta z_n$ relative to $z_n^\star$.
The \code{SubType} parameter selects this lower-precision type.  It
is typically \code{float}, \code{double}, or
\code{CudaDblflt<MattDblflt>}, and for HDR-wrapped references the
corresponding \code{HDRFloat} variant.  The mapping from \code{T} to
\code{SubType} is encoded by the \code{SubTypeChooser} trait: when
\code{T} is a plain scalar the sub-type is the same scalar; when
\code{T} is an \code{HDRFloat<S>} the sub-type extracts the
underlying~\code{S}.

\code{SubType} appears as an explicit template parameter on
\code{LAReference<IterType,\allowbreak T,\allowbreak SubType,\allowbreak
PExtras>} and is threaded into the perturbation kernels, where it
controls the precision of the per-pixel inner loop.


\subsection{Perturbation extras (\code{PExtras})}
\label{subsec:template-pextras}

The \code{PerturbExtras} enumeration is a compile-time flag attached to
every \code{PerturbationResults} instantiation.  It controls whether the
orbit stores auxiliary data beyond the orbit values themselves:

\begin{itemize}
  \item \code{Disable}---store only the orbit; no extra bookkeeping.
  \item \code{Bad}---record \emph{glitch} information (iterations where
        the perturbation approximation loses accuracy) alongside the
        orbit, enabling downstream glitch-correction passes.
  \item \code{SimpleCompression} and \code{MaxCompression}---store the
        metadata required for reference-orbit compression
        (\cref{sec:ref-compression-zhuoran}).
\end{itemize}

\noindent
Making this a template parameter rather than a run-time flag eliminates
the storage overhead of the auxiliary arrays in the common case where
they are not needed, and allows the compiler to elide the associated
bookkeeping code entirely.


\subsection{Instantiation combinatorics}
\label{subsec:template-combinatorics}

The Cartesian product of the four dimensions yields a large number of
potential type combinations.  With two \code{IterType} widths, six
representative \code{T} values, and three commonly used
\code{PExtras} settings (\code{Disable}, \code{Bad},
\code{Simple\-Compression}), the \code{PerturbationResults} class
alone accounts for $2\times 6\times 3 = 36$ concrete instantiations.
This count is visible in the \code{AwesomeVariant} type alias defined
in \code{RefOrbitCalc.h}, which enumerates all 36~alternatives as a
\code{std::variant} so that perturbation results of any type
combination can be stored and dispatched uniformly at run time.

Additional template parameters (\code{SubType}, benchmark modes, reuse
modes) on related classes further multiply the instantiation count.  The
full set of explicit instantiations in
\code{RefOrbitCalcTemplates.h} runs to several hundred
macro-expanded declarations.


\subsection{Explicit instantiation strategy}
\label{subsec:template-explicit-inst}

Because every translation unit that includes a template header would
otherwise independently instantiate each specialisation it references,
\FractalShark{} relies on \emph{explicit template instantiation} to
centralise code generation and control compile times.

For the perturbation pipeline, the header
\code{RefOrbitCalcTemplates.h} defines a family of preprocessor macros
(\code{InstantiateIsPerturbation\-Result\-UsefulHere},
\code{InstantiateAddPerturbation\-Reference\-Point}, and others) that
expand to explicit instantiation declarations for every supported
\code{(IterType, T, SubType, PExtras)} combination.  Each macro accepts
the type arguments, expands to the required \code{template} declaration,
and is invoked once per combination immediately below its definition.
This pattern confines the instantiation cost to a single translation
unit while allowing other files to use only the declarations.

For the GPU arithmetic subsystem, a separate code-generation tool
(\code{HpShark\-Instantiate}) produces \code{.cu} files containing
explicit instantiations of \code{HpSharkFloat}, the NTT multiply
kernels, the reference-orbit GPU loop, and associated helper templates.
The generator partitions the output into paired
files---one for production parameter sets
(\code{SharkParams1}\,--\,\code{12}),
one for non-periodicity variants
(\code{SharkParamsNP1}\,--\,\code{NP12})---and groups
instantiations into named \emph{batches} (e.g.\
\code{HpSharkFloat\_Conversions}, \code{MultiplyNTT},
\code{HpSharkReference}) so that each batch compiles as an
independent translation unit.  This decomposition enables parallel
compilation and limits the peak memory consumption of the CUDA
compiler.


\subsection{The \code{SharkFloatParams} pattern}
\label{subsec:template-sharkparams}

The GPU high-precision arithmetic library takes the compile-time
parameterisation one step further by encoding not just type choices but
also \emph{algorithmic configuration} in a template parameter.
The struct
\code{GenericSharkFloatParams<pNumDigits,\allowbreak Periodicity,\allowbreak
NewtonRaphson,\allowbreak SubTypeT>}
bundles several \code{static constexpr} members that the CUDA compiler
can fold into kernel code:

\begin{itemize}
  \item \code{GlobalNumUint32}---the number of 32-bit limbs in each
        high-precision operand, derived from \code{pNumDigits}.  Because
        the compiler knows this value at template instantiation time,
        loop trip counts in the NTT multiply and carry-propagation
        kernels are compile-time constants, enabling aggressive
        unrolling and register allocation.
  \item \code{EnablePeriodicity}---controls whether the reference-orbit
        kernel checks for period detection at each iteration.  Disabling
        it in the non-periodicity (\code{SharkParamsNP*}) and
        Newton--Raphson (\code{SharkParamsNR*}) variants eliminates
        the associated branch and storage overhead.
  \item \code{EnableNewtonRaphson}---when set, the kernel additionally
        tracks the derivative $d_n = \partial z_n/\partial c$ alongside
        the orbit, allocating extra GPU scratch memory and executing the
        derivative recurrence at each iteration.  This is required by
        the Feature Finder's Newton's method solver
        (\cref{sec:feature-finder}).
  \item \code{NTTPlan}---a compile-time NTT plan computed by
        \code{SharkNTT::BuildPlanPrime} from \code{GlobalNumUint32},
        fixing the NTT transform length and modular-arithmetic
        parameters at compile time.
\end{itemize}

\noindent
Concrete aliases such as
\code{SharkParams1}\,--\,\code{SharkParams12} fix
\code{pNumDigits} to powers of two from $256$ through $524{,}288$
(providing roughly $2{,}500$ to $5{,}000{,}000$ decimal digits of
precision) with periodicity enabled and Newton--Raphson disabled.
Additional alias families---\code{SharkParamsNP*} (no periodicity),
\code{SharkParamsNR*} (Newton--Raphson enabled),
\code{SharkParamsDbl*} (\code{double} sub-type), and
\code{SharkParamsDbf*} (\code{CudaDblflt} sub-type)---provide the
remaining feature-flag and sub-type permutations.  Each alias is
independently instantiated by the code generator described in
\cref{subsec:template-explicit-inst}, yielding $12\times 5 = 60$
parameter-set instantiations per GPU kernel template.

With this architectural overview in place, the following chapter turns
to the rendering kernels themselves, beginning with the pixel-to-parameter
mapping and the base escape-time iteration that all subsequent optimisations
build upon.
